\section{Escopo}

\paragraph{}
O escopo deste teste de penetração foi definido com base nas informações
obtidas durante a varredura de vulnerabilidades realizada com a ferramenta
OWASP ZAP no ambiente disponibilizado pelo Laboratório de Segurança.

\subsection{Itens incluídos no escopo}

\begin{itemize}
\item \textbf{Aplicação Web Principal}
\begin{itemize}
\item URL: \texttt{http://testasp.vulnweb.com}
\item Tipo: Aplicação web pública de demonstração.
\item Tecnologias observadas: ASP clássico e servidor web Windows.
\item Objetivo: Avaliação de vulnerabilidades em funcionalidades, parâmetros,
  formulários, mecanismos de autenticação e controle de acesso.
\end{itemize}

\item \textbf{Páginas e funcionalidades identificadas durante o mapeamento}
\begin{itemize}
    \item \texttt{/Search.asp}
    \item \texttt{/showthread.asp}
    \item \texttt{/Templatize.asp}
    \item \texttt{/Login.asp}
    \item Demais rotas descobertas durante o processo de spider e varredura
      automatizada.
\end{itemize}

\item \textbf{Parâmetros vulneráveis identificados}
\begin{itemize}
    \item \texttt{tfSearch} (SQL Injection);
    \item \texttt{tfText} (Cross Site Scripting Refletido);
    \item \texttt{item} (Path Traversal);
    \item \texttt{RetURL} (External Redirect).
\end{itemize}

\item \textbf{Controles de segurança da aplicação}
\begin{itemize}
    \item Proteções contra CSRF;
    \item Cabeçalhos HTTP de segurança;
    \item Configurações de cookies;
    \item Políticas de Content Security Policy (CSP);
    \item Proteções contra Clickjacking.
\end{itemize}

\end{itemize}

\subsection{Áreas prioritárias do teste}

As principais áreas de interesse do Pentest são aquelas que apresentaram
vulnerabilidades classificadas como de alto risco durante a análise com o OWASP
ZAP:

\begin{itemize}
\item SQL Injection (CWE-89);
\item Cross Site Scripting Refletido -- XSS (CWE-79);
\item Path Traversal (CWE-22);
\item External Redirect (CWE-601).
\end{itemize}

Essas vulnerabilidades possuem potencial para comprometer a confidencialidade,
integridade e disponibilidade das informações da aplicação, sendo consideradas
prioritárias durante a execução dos testes.

\subsection{Itens excluídos do escopo}

\begin{itemize}
\item Sistemas externos ao domínio \texttt{testasp.vulnweb.com};
\item Infraestruturas de terceiros não relacionadas ao laboratório;
\item Equipamentos pessoais dos participantes;
\item Redes externas à aplicação avaliada;
\item Ataques de negação de serviço (DoS/DDoS);
\item Engenharia social;
\item Qualquer sistema não explicitamente autorizado pela atividade.
\end{itemize}

\paragraph{}
Dessa forma, o escopo concentra-se na avaliação de segurança da aplicação web e
de seus componentes diretamente relacionados, priorizando as áreas em que a
varredura inicial identificou vulnerabilidades de maior criticidade.
